% !TEX root = ../main.tex

\chapter{Conclusions}

\section{Discussions}

The central question of this thesis is whether a mobile literature assistant can preserve continuity across changing papers, generated evidence, and changing user interests.
The current evidence supports a qualified answer.
The system design and inspected implementation establish mechanisms for preserving identity, evidence scope, event semantics, and device-visible state at the stated cutoff.
They do not yet establish final quality across a representative document corpus, question set, longitudinal behavior set, and device study.
The discussion therefore answers each research question by separating what the mechanism establishes from what its evaluation must still determine.

For revision-correct representation, the design and inspected implementation separate logical paper identity from revision identity and carry revision, content hash, transformation configuration, and lineage into derived artifacts.
Revision-scoped chunks and the retrieval-isolation test establish that the mechanism is executable and that the tested Q\&A boundary does not mix revisions.
The retained evidence does not establish parser structure quality, invalidation completeness, recovery rates, or performance across a diverse scientific-document corpus, so these properties remain outside the conclusion.

For evidence-bounded assistance, the design decomposes answers into claims and distinguishes identifier validity, evidence availability, consistency, and coverage.
It defines scoped retrieval, resolvable source spans, and abstention.
The literature supports the need for these distinctions~\cite{gao2023alce,es2024ragas,huang2024grounded,niu2024ragtruth}.
The retained software evidence establishes the presence of these controls but does not establish a general improvement in retrieval or citation quality.

For adaptive interest memory, the design combines explicit interests, bounded implicit-event semantics, a time-decayed normalized behavioral representation, inspectable recommendation reasons, and user correction.
The available artifacts do not establish improvements in ranking, diversity, longitudinal adaptation, or user control, and the conclusion is therefore limited to the implemented and inspected mechanism.

For bounded mobile integration, the Android implementation and prototype demonstrate that the intended workflow can be represented as mobile interactions with local cache, recoverable event delivery, and a bounded graph boundary.
The executed Android subset provides measured software evidence on one Android~14 emulator: all retained startup/resume trials, controlled state transitions, graph ready/selection samples, and controlled event stages met their defined success conditions, while a five-run live trace confirmed duplicate-aware event ingestion and client-loop reachability.
The supplied presentation reports that participants completed the digest, evidence, Q\&A, and preference tasks more often than the graph task.
Because participant-level records and a physical-device matrix are absent, the usability evidence is \observed{}/\limited{} and the measured device result is restricted to the stated emulator/configuration.
Notification benefit, privacy behavior, recommendation adaptation, and graph-redesign usefulness remain unestablished.

Taken together, the contributions and limitations define a deliberately narrow claim envelope.
The first contribution is a revision-aware data and job model that makes paper-derived state attributable and recoverable.
The second is an evidence-bounded assistance contract that distinguishes retrieved evidence, marker validity, limited source matching, and refusal.
The third is a versioned behavioral contract and auditable recommendation baseline whose reasons derive from the same signals used in scoring.
The fourth is a mobile integration design that separates live, cached, fixture, fallback, and error states while returning bounded events to the server.
The fifth is an evaluation framework that maps these contributions to revision, evidence, adaptation, and mobile-system measurements and prevents software tests from being reported as research effectiveness.

The same evidence-first standard limits the claim envelope.
The ten customer interviews and five prototype sessions are formative, small, and concentrated in computing-related contexts.
Full transcripts, consent artifacts, participant-level usability tables, and final retest records are unavailable in the local package.
The 15-case Q\&A fixture is an engineering seed rather than a benchmark.
The current source-matching verifier measures lexical overlap, not semantic entailment or scientific truth.
Behavior-v1 uses hand-specified weights and ambiguous proxies rather than longitudinally validated intent.
No retained evidence supports a claim about representative document diversity, answer quality at benchmark scale, longitudinal recommendation benefit, a physical-device matrix, or post-redesign participant performance.
Consequently, the thesis supports an implementation-grounded design contribution and bounded mobile measurements, but not a general claim that MNEME improves research productivity or understanding.

\section{Main Conclusions}

This thesis began from a downstream literature problem: after a paper has been found, researchers must still judge it, remember it, reconnect it to later work, and verify machine-produced explanations.
MNEME addresses this problem as an accountable mobile loop rather than as an autonomous scientist.
The loop receives and prepares changing papers, produces source-linked assistance, records explicit and bounded behavioral signals, assembles a versioned briefing, and exposes a limited citation neighborhood on Android.

The main technical result at the present cutoff is the alignment of four contracts.
Revision identity and artifact lineage protect document continuity.
Scoped retrieval, explicit source items, refusal, and limited verification protect evidence continuity.
Versioned events, an auditable scorer, and editable interests protect preference continuity.
Room persistence, WorkManager recovery, bounded graph state, and visible content origins carry those properties into the mobile workflow.
The reported formative study indicates that digest, source, and cited-Q\&A paths were discoverable for the five participants, while graph interaction required a clearer next action.

The evidence does not justify a uniformly positive answer to all four research questions.
It establishes that the mechanisms can be specified, implemented, and tested at software level, and that the retained Android path operates within the reported emulator and live-service boundaries.
MNEME is therefore most defensible as an implementation-grounded research system with bounded mobile evidence.
Its central lesson is that an AI literature assistant should not be judged by fluent output or feature count alone: its paper versions must remain identifiable, its generated claims inspectable, its behavioral adaptation correctable, and its failures visible.

\section{Outlook}

The first priority is an immutable final repository, build, configuration, and corpus freeze.
A broader revision study should then measure paper identity, parser quality, stale-artifact invalidation, recovery, and latency on a stratified scientific-document set.
An evidence-quality study should execute a frozen question suite, add human claim--evidence judgments, compare retrieval and verification ablations, and report refusal, latency, and cost separately.
A longitudinal adaptation study should collect exposure-aware temporal event sequences, compare explicit-only, recency-based, and behavioral models, and test stability, diversity, and user correction.
A mobile follow-up should complete live briefing refresh and notification behavior, execute a physical-device reliability matrix, retain participant-level task records, and repeat the graph task on the revised connected build.

A second research phase should broaden disciplines, document layouts, revision patterns, and longitudinal duration.
Privacy work should test deletion, correction, export, and retention behavior rather than treating data minimization as a design statement.
Semantic support checking may replace or augment lexical source matching only after calibration against human labels.
Graph work can examine concept and author layers after the bounded citation interaction is validated.

Stretch features such as autonomous multi-hop research, voice input, camera OCR, shared team lists, reference-manager export, and on-device language models should follow the core evidence chain.
They broaden capability but do not repair missing evidence for provenance, support, adaptation, or mobile reliability.
The project should therefore prioritize claim closure before feature expansion.

\clearpage
